\chapter{Applications to many-body systems}
\label{chap:appmb}
%=============================================================================

Every method in this book has now been developed, and every one of them has
been tested on a system chosen to make it easy.  Hartree-Fock, perturbation
theory and coupled cluster were applied in Chapter~\ref{chap:applications} to a
quantum dot in a finite oscillator basis, where the two-body matrix elements
are analytic.  Variational and diffusion Monte Carlo were applied in
Chapters~\ref{chap:vmc} to \ref{chap:dmc} to \emph{two} electrons in the same
trap, where the trial function is nodeless and a determinant is not needed at
all.

This chapter puts all five on the same problem: six electrons in a
two-dimensional parabolic trap.  Six is the smallest number for which none of
the simplifications survive.  The determinant is no longer a formality, the
trial function has nodes, the basis-set methods are no longer converged, and
there is no exact answer to check against -- only the internal consistency of
five methods that fail in different ways, and published benchmarks.

Two things have to be built.  The first is the Slater-Jastrow trial function
and the linear algebra that makes it evaluable, which is the material of
Sections~\ref{sec:slaterupdate} to \ref{sec:slaterstability}: the determinant
ratio from the inverse Slater matrix, the gradient and Laplacian ratios that
supply the quantum force and the kinetic energy, and the factorisation
$\Det{D}=\Det{D}_\uparrow\Det{D}_\downarrow$.  That material was written in
Chapter~\ref{chap:mbphys} without a physical application; this is the
application.

The second is the spin.  Everything in Chapters~\ref{chap:vmc} to
\ref{chap:dmc} sampled positions.  Here the assignment of electrons to the two
spin blocks is itself a degree of freedom, and we sample it -- which turns the
Moskowitz-Kalos argument of Section~\ref{sec:spinfactorisation} from an
assertion into a measurement, and which behaves quite differently in the two
Monte Carlo methods.  That difference is the most interesting thing in the
chapter, and it is a trap.

The programs are in \texttt{manybodymc.py}, which imports the oscillator
orbitals and their derivatives from \texttt{slater\_update.py}, the blocking
analysis from \texttt{vmcoptimise.py}, and reproduces \texttt{vmc.py} exactly
when $N=2$.

%----------------------------------------------------------------------
\section{The system, and what is already known about it}
\label{sec:appsystem}

\index{quantum dot!six electrons}
The Hamiltonian is the one of Section~\ref{sec:qdsystem}, in the atomic units
used there,
\begin{equation}
  \hat{H} = \sum_{i=1}^{N}\left(-\frac{1}{2}\nabla_i^2
    + \frac{1}{2}\omega^2 r_i^2\right)
    + \sum_{i<j}\frac{1}{r_{ij}} ,
  \label{eq:16-hamiltonian}
\end{equation}
with $N=6$ and $\hbar\omega=1$ unless stated otherwise.  Six electrons fill the
two lowest oscillator shells -- one orbital in the first, two in the second,
each doubly occupied -- so the ground state is a closed shell with $S=0$ and
three electrons of each spin.  This is what made $N=6$ one of the magic numbers
of Section~\ref{sec:qdbasis}.

Chapter~\ref{chap:applications} has already solved it three ways, in a basis of
six oscillator shells:
\begin{equation}
  E_{\mathrm{HF}} = 20.72026, \quad
  E_{\mathrm{MP2}} = 20.30256, \quad
  E_{\mathrm{CCSD}} = 20.27325 .
  \label{eq:16-basisresults}
\end{equation}
None of these is converged.  The trouble is the cusp: the exact wave function
has a kink at $r_{ij}=0$, no finite sum of products of oscillator orbitals can
reproduce a kink, and the coefficients therefore decay slowly with the number
of shells.  Lohne \emph{et al.}~\cite{lohne2011} carried the same CCSD
calculation to twenty shells and obtained $20.1737$; the $42$-orbital number
above is a tenth of an atomic unit above it.  Whatever else this chapter shows,
it is not a fair contest between coupled cluster and Monte Carlo -- it is a
demonstration that a basis-set calculation must be converged before it is
compared with anything.

The benchmark is the fixed-node diffusion Monte Carlo result of the same
reference,
\begin{equation}
  E_{\mathrm{DMC}} = 20.1597(2),
  \label{eq:16-benchmark}
\end{equation}
against which their best coupled-cluster number, $\Lambda$-CCSD(T) in twenty
shells, gives $20.1582$.  Those two agreeing to $1.5\times10^{-3}$ from
opposite directions is the strongest statement available about this system, and
it is what our own calculation has to reproduce.

%----------------------------------------------------------------------
\section{The Slater-Jastrow trial function}
\label{sec:appansatz}

\index{Slater-Jastrow wave function}
The trial function is
\begin{equation}
  \boxed{\;
  \Psi_T(\bm{R},\bm{\sigma})
   = \Det{D}_{\uparrow}\,\Det{D}_{\downarrow}\;
     \exp\left(\sum_{i<j}\frac{a_{ij}\,r_{ij}}{1+\beta r_{ij}}\right) \;}
  \label{eq:16-slaterjastrow}
\end{equation}
with $\Det{D}_{\uparrow}$ and $\Det{D}_{\downarrow}$ the $3\times3$
determinants of Eq.~(\ref{eq:2-detfactorises}) over the spin-up and spin-down
electrons, built from the oscillator orbitals
$\phi_{n_xn_y}(\bm{r};\alpha\omega)$ at the \emph{scaled} frequency
$\alpha\omega$.  There are exactly two variational parameters, $\alpha$ and
$\beta$, as in Chapter~\ref{chap:vmc}.

For $N=2$ both determinants are $1\times1$ and equal to the oscillator ground
state, and Eq.~(\ref{eq:16-slaterjastrow}) collapses onto the Pad\'e-Jastrow
function of Eq.~(\ref{eq:13-trial}).  That is worth checking numerically rather
than by inspection, and \texttt{manybodymc.py} does: for the same
configuration and the same parameters, its local energy agrees with that of
\texttt{vmc.py} to twelve digits.

\paragraph{The cusp coefficients are spin dependent.}
\index{cusp condition!spin dependence}
The $a_{ij}$ are not variational.  They are fixed by requiring $E_L$ to stay
finite as two electrons meet, and the requirement is \emph{not} the same for
the two spin cases.  Repeating the argument of Section~\ref{sec:vmctrial} in
$d$ dimensions for a pair whose relative wave function behaves as $r^{\ell}$ at
contact gives
\begin{equation}
  a = \frac{1}{2\ell + d - 1} .
  \label{eq:16-cusp}
\end{equation}
For an antiparallel pair the spatial function is finite at contact, $\ell=0$,
and $a=1$ in two dimensions.  For a parallel pair the two electrons sit in the
same determinant, which vanishes linearly when they coincide, so $\ell=1$ and
\begin{equation}
  a_{\uparrow\uparrow} = a_{\downarrow\downarrow} = \frac{1}{3},
  \qquad
  a_{\uparrow\downarrow} = 1 .
  \label{eq:16-cuspvalues}
\end{equation}
The determinant has already done a third of the work, and the Jastrow factor
must not do it again.

Table~\ref{tab:16-cusp} shows what happens if this is forgotten.  Two electrons
are brought together in a six-electron configuration, once as an antiparallel
pair and once as a parallel one, and the third column repeats the parallel case
with $a_{\uparrow\uparrow}=0$ -- the mistake of assuming that the determinant
handles the parallel case by itself.

\begin{table}[h]
\centering
\renewcommand{\arraystretch}{1.2}
\setlength{\tabcolsep}{14pt}
\begin{tabular}{llll}
\toprule
$r_{ij}$ & antiparallel, $a=1$ & parallel, $a=\tfrac13$
  & parallel, $a=0$\\
\midrule
$1.000$ & $17.0116$ & $16.7744$ & $17.03$\\
$0.300$ & $14.1059$ & $15.1091$ & $17.82$\\
$0.100$ & $12.3815$ & $14.4457$ & $23.88$\\
$0.030$ & $11.4901$ & $14.3025$ & $47.09$\\
$0.010$ & $11.1825$ & $14.2804$ & $113.74$\\
$0.003$ & $11.0676$ & $14.2751$ & $347.07$\\
\bottomrule
\end{tabular}
\caption{The local energy as two of six electrons are brought together, with
the other four held fixed.  With the correct coefficients the local energy
converges; with $a_{\uparrow\uparrow}=0$ it diverges as $1/r_{ij}$, and the
variance of the energy is then unbounded.  Computed by
\texttt{manybodymc.py}.}
\label{tab:16-cusp}
\end{table}

\begin{notebox}
\textbf{Why this is the interesting parameter.}  The spin dependence of
$a_{ij}$ is what makes the spin assignment a physical variable rather than
bookkeeping.  If the Jastrow factor were spin blind, exchanging the spins of
two electrons would change $\Psi_T$ only through the determinants; with
Eq.~(\ref{eq:16-cuspvalues}) it changes the correlation factor with every other
electron as well.  That is why the exchange move of
Section~\ref{sec:appspin} has an acceptance rate worth measuring.
\end{notebox}

%----------------------------------------------------------------------
\section{The determinant machinery}
\label{sec:appdeterminant}

\index{Slater determinant!efficient evaluation}
A Monte Carlo run makes millions of single-particle moves and needs, at each
one, the ratio $\Psi_T(\bm{R}')/\Psi_T(\bm{R})$ and then -- if the move is
accepted -- the gradient and Laplacian of $\ln\Psi_T$ for every particle.
Section~\ref{sec:slaterupdate} showed that all of this follows from the
\emph{inverse} Slater matrix.  Nothing needs to be rederived here; it needs to
be used.

\paragraph{The ratio.}
Moving particle $i$ changes one row of one of the two blocks, so by
Eq.~(\ref{eq:2-Rinverse})
\begin{equation}
  R = \sum_{j}\phi_j(\bm{r}_i^{\mathrm{new}})\,d^{-1}_{ji}
  \label{eq:16-ratio}
\end{equation}
is a single dot product of length $N/2$.  The other block is untouched and
cancels, which is the saving of Eq.~(\ref{eq:2-ratiofactorises}).  If the move
is accepted the inverse is repaired by the Sherman-Morrison
formulae~(\ref{eq:2-inverseupdate}).  Two hundred consecutive accepted moves in
\texttt{manybodymc.py} reproduce the brute-force quotient of two determinants
to $3\times10^{-12}$, and $\|\bm{D}^{-1}\bm{D}-\bm{1}\|$ is still
$9\times10^{-15}$ at the end: the inverse does not drift and never has to be
rebuilt.

\paragraph{The derivatives.}
Equations~(\ref{eq:2-gradratio}) and (\ref{eq:2-lapratio}) give
$\nabla_i\Det{D}/\Det{D}$ and $\nabla_i^2\Det{D}/\Det{D}$ as the same kind of
dot product, with $\phi_j$ replaced by its derivatives.  The local energy then
follows from
\begin{equation}
  \frac{\nabla^2\Psi_T}{\Psi_T}
   = \frac{\nabla^2\Det{D}}{\Det{D}}
   + 2\,\frac{\nabla\Det{D}}{\Det{D}}\cdot\nabla U
   + \nabla^2 U + \left|\nabla U\right|^2 ,
  \label{eq:16-laplacian}
\end{equation}
$U$ being the exponent of the Jastrow factor.  The determinant and the
correlation factor are never differentiated together: only their separate
derivatives appear, and each is available in closed form.  The quantum force is
\begin{equation}
  \bm{F}_i = 2\left(\frac{\nabla_i\Det{D}}{\Det{D}} + \nabla_i U\right).
  \label{eq:16-quantumforce}
\end{equation}
Every one of these is differentiated by hand in \texttt{manybodymc.py} and
checked against finite differences before anything is sampled -- the gradient
of $\ln\Psi_T$ to $4\times10^{-8}$ and the local energy to $1\times10^{-4}$,
both at the level of the finite-difference truncation error rather than of the
formulae.

\begin{notebox}
\textbf{A practical note on the scaling.}  Table~\ref{tab:vmcscaling} promises
$\bigO(dN^2)$ instead of $\bigO(dN^4)$, and that is the reason the algorithm is
written this way.  For $N=6$, however, the determinants are $3\times3$, and a
vectorised batched factorisation of a thousand such matrices at once is faster
in practice than a Python loop carrying Sherman-Morrison updates one walker at
a time.  \texttt{manybodymc.py} therefore contains both: the single-walker
class that states the algorithm and verifies it, and a batched evaluator used
for the production runs.  The asymptotic argument is right and the crossover is
simply well above $N=6$.
\end{notebox}

%----------------------------------------------------------------------
\section{Sampling the spins}
\label{sec:appspin}

\index{spin sampling}\index{Moskowitz-Kalos ansatz}
Section~\ref{sec:spinfactorisation} argued that for a spin-independent
Hamiltonian the full antisymmetric determinant may be replaced by the product
$\Det{D}_{\uparrow}\Det{D}_{\downarrow}$ at a fixed assignment of electrons to
the two blocks, and that the energy is unchanged.  Every production code takes
that on trust and never touches the assignment again.

We do not have to.  Treat the spin configuration
$\bm{\sigma}=(\sigma_1,\dots,\sigma_N)$, $\sigma_i=\pm1$, as a coordinate
alongside the positions, and add a move: pick one spin-up and one spin-down
electron at random and exchange their spins.  The proposal is symmetric -- the
reverse move draws the same pair with the same probability -- so the ordinary
Metropolis test applies,
\begin{equation}
  q = \frac{\left|\Psi_T(\bm{R},\bm{\sigma}')\right|^2}
           {\left|\Psi_T(\bm{R},\bm{\sigma})\right|^2} ,
  \label{eq:16-spinmetropolis}
\end{equation}
and $S_z$ is conserved exactly.  The move is not trivial: the two electrons
change determinants, and by Eq.~(\ref{eq:16-cuspvalues}) their cusp
coefficients with every other electron change too.  For six electrons at the
optimised parameters it is accepted $42\%$ of the time.

What the move buys is discussed in the next section, and it is not what one
would expect.  What it settles immediately is the Moskowitz-Kalos claim: with
the spins sampled the variational energy is $20.2000(18)$, and with the
assignment frozen it is $20.1956(31)$.  The two agree, as promised, and now as
a measurement rather than an argument.

%----------------------------------------------------------------------
\section{Nodes, and two ways to survive them}
\label{sec:appnodes}

\index{nodal surface}\index{drift!divergence at nodes}
Here is the first thing that breaks in going from two electrons to six.  The
two-electron trial function of Chapter~\ref{chap:vmc} is strictly positive.
This one is not: $\Det{D}_{\uparrow}$ vanishes whenever the three spin-up
electrons become collinear, which for the orbitals $(0,0)$, $(1,0)$, $(0,1)$ is
exactly the condition that
\begin{equation}
  \det\begin{pmatrix}
    1 & x_1 & y_1\\ 1 & x_2 & y_2\\ 1 & x_3 & y_3
  \end{pmatrix} = 0 .
  \label{eq:16-node}
\end{equation}
A walker can approach that surface, and when it does the quantum force
$\bm{F}=2\nabla\ln\Psi_T$ diverges.  The drift step $D\bm{F}\Delta t$ then
hurls the walker across the configuration space, and the local energy it
reports on the way is enormous.  In a run of three hundred walkers the
walker-averaged energy swings between $18.6$ and $21.7$ and the estimator is
worthless -- not because of any error in the formulae, all of which have been
checked, but because the Langevin proposal of Eq.~(\ref{eq:13-proposal}) is
inappropriate near a singular drift.

There are two cures, and it is instructive that they are independent.

\paragraph{Cut the drift.}
\index{effective time step}
Umrigar, Nightingale and Runge~\cite{umrigar1993} replace the drift by
\begin{equation}
  \bm{F} \;\longrightarrow\;
  \bm{F}\,\frac{-1+\sqrt{1+2a\bm{F}^2\Delta t}}{a\bm{F}^2\Delta t},
  \label{eq:16-driftcutoff}
\end{equation}
which is the identity for small $\bm{F}$ and saturates for large, so the drift
step never exceeds a fixed size.  The essential point is that the \emph{same}
modified drift is used in the proposal and in the Green's function of
Eq.~(\ref{eq:13-greensratio}).  Detailed balance is then untouched and the
sampled distribution is still exactly $|\Psi_T|^2$ -- which is the lesson of
Section~\ref{sec:vmcimportance}, where a deliberately wrong quantum force was
shown to cost efficiency and not correctness.  Here that freedom is used on
purpose.

\paragraph{Or let the spins move.}
A walker that has strayed close to a node has a small $|\Psi_T|$, so by
Eq.~(\ref{eq:16-spinmetropolis}) almost any spin exchange is accepted:
exchanging one of the three collinear electrons into the other block destroys
the near-degeneracy at once.  The exchange move is an escape hatch that opens
precisely when it is needed, and it requires the walker to move no distance at
all.

Table~\ref{tab:16-nodes} runs all four combinations.

\begin{table}[h]
\centering
\renewcommand{\arraystretch}{1.2}
\setlength{\tabcolsep}{11pt}
\begin{tabular}{llll}
\toprule
drift & spin moves & energy & series range\\
\midrule
plain & no  & $20.2547(498)$ & $18.547$ -- $21.439$\\
plain & yes & $20.2036(18)$  & $20.127$ -- $20.260$\\
cut   & no  & $20.1956(31)$  & $20.108$ -- $20.249$\\
cut   & yes & $20.2000(18)$  & $20.135$ -- $20.261$\\
\bottomrule
\end{tabular}
\caption{Variational Monte Carlo for six electrons at $\alpha=1.00$,
$\beta=0.46$, $300$ walkers and $400$ steps.  ``Series range'' is the smallest
and largest walker-averaged energy over the run.  Computed by
\texttt{manybodymc.py}.}
\label{tab:16-nodes}
\end{table}

The three treated rows agree with each other; the untreated one does not, and
its error bar does not warn you.  Two further observations are worth making.
The spin exchange alone is enough -- the second row needs no drift cutoff at
all -- which says that the pathology really is walkers dwelling near nodes
rather than anything in the drift formula.  And where both are available the
exchange is still the better of the two: its error bar is smaller by a factor
between $1.3$ and $1.7$ across several seeds, which is a factor of two or so in
computer time for the same accuracy.  Both are used from here on.

%----------------------------------------------------------------------
\section{Optimising the trial function}
\label{sec:appoptimise}

The machinery of Chapter~\ref{chap:optimres} carries over untouched.  The
energy gradient is still the covariance of Eq.~(\ref{eq:14-gradient}), and only
the parameter derivatives change.  That with respect to $\beta$ is elementary.
That with respect to $\alpha$ now runs through the orbitals of the determinant,
and is assembled blockwise as
\begin{equation}
  \frac{\partial \ln\Det{D}}{\partial \alpha}
   = \mathrm{Tr}\left(\bm{D}^{-1}\frac{\partial \bm{D}}{\partial\alpha}\right),
  \label{eq:16-alphagradient}
\end{equation}
which is the same inverse the sampling already maintains.  Both derivatives
agree with finite differences of $\ln\Psi_T$ to $1\times10^{-9}$.

Started deliberately far from the minimum, at $\alpha=0.90$ and $\beta=0.25$,
plain gradient descent with $150$ walkers and $120$ steps per iteration --
short, noisy runs, as Chapter~\ref{chap:optimres} insisted -- moves
monotonically to $\alpha=1.00$, $\beta=0.40$ in fourteen iterations, lowering
the energy from $21.008$ to $20.221$.  Refining on a scan puts the minimum at
\begin{equation}
  \alpha = 0.98, \qquad \beta = 0.46 ,
  \label{eq:16-optimum}
\end{equation}
with the surface flat to within $8\times10^{-3}$ over
$\alpha\in[0.96,1.02]$ and essentially flat in $\beta$ between $0.44$ and
$0.48$.

\begin{svgraybox}
That $\alpha$ comes out at $1$ to within the noise is the physically
interesting result.  It says that the optimal orbitals in the determinant are
the \emph{non-interacting} oscillator eigenfunctions: the Jastrow factor
absorbs the effect of the Coulomb interaction on the one-body part so
completely that there is nothing left for a scale parameter to do.  Harju and
co-workers~\cite{harju2002} found the same for this system by optimising the
single-particle states freely, and noted that it fails for the three-electron
dot, where the numbers of spin-up and spin-down electrons differ.  Closed
shells are forgiving.
\end{svgraybox}

%----------------------------------------------------------------------
\section{Variational Monte Carlo}
\label{sec:appvmc}

At the optimum, six hundred walkers taking nine hundred steps -- $540\,000$
samples, in about half a minute -- give
\begin{equation}
  E_{\mathrm{VMC}} = 20.19961 \pm 0.00091 ,
  \label{eq:16-vmc}
\end{equation}
with an acceptance rate of $97.4\%$, $41.6\%$ of the spin exchanges accepted,
and a correlation time of $2.74$ steps.  The error is from blocking; the
bootstrap gives $0.00065$ on the same series, and the two agree on the order of
magnitude in the way Section~\ref{sec:optresampling} led one to expect.

This single number already sits $0.074$ below the CCSD energy of
Eq.~(\ref{eq:16-basisresults}), and $0.040$ above the benchmark of
Eq.~(\ref{eq:16-benchmark}).  Neither comparison is quite what it looks like.
The first is a statement about the oscillator basis, not about coupled cluster.
The second is the deficiency of a two-parameter trial function, and it is what
diffusion Monte Carlo now removes.

%----------------------------------------------------------------------
\section{Diffusion Monte Carlo, and a warning}
\label{sec:appdmc}

\index{diffusion Monte Carlo!fixed node}
The algorithm is that of Chapter~\ref{chap:dmc}, unchanged except that the
guiding function is now Eq.~(\ref{eq:16-slaterjastrow}) and that the fixed-node
approximation of Section~\ref{sec:dmcfixednode} is no longer vacuous.  The
nodal surface is Eq.~(\ref{eq:16-node}), and one thing about it is worth
noticing: the Gaussian factors are common to all rows and cancel from the
determinant, so the node depends on $\alpha$ only through a uniform dilation
of the configuration space.  For six electrons the nodal surface is therefore
essentially fixed by symmetry, and the fixed-node energy is a well-defined
number rather than something that drifts as the trial function is retuned.

\paragraph{The spin move must now be switched off.}
This is the trap.  The exchange move of Section~\ref{sec:appspin} satisfies
detailed balance with respect to $|\Psi_T|^2$.  That is exactly what
variational Monte Carlo samples, and it is exactly \emph{not} what diffusion
Monte Carlo samples: the target there is $f=\Psi_T\Phi$ of
Eq.~(\ref{eq:15-mixeddistribution}).  Applying the exchange at every step
inserts a move that pulls the walkers back towards the variational
distribution, in competition with the branching that is trying to project them
away from it.  The result is not a small bias.

\begin{table}[h]
\centering
\renewcommand{\arraystretch}{1.2}
\setlength{\tabcolsep}{14pt}
\begin{tabular}{ll}
\toprule
spin exchanges & energy\\
\midrule
at every step & $20.19483(321)$\\
during equilibration only & $20.14812(507)$\\
never & $20.14980(340)$\\
\bottomrule
\end{tabular}
\caption{Diffusion Monte Carlo at $\Delta\tau=0.02$ with $200$ walkers.
Exchanging the spins throughout returns, to within its error bar, the
variational energy of Eq.~(\ref{eq:16-vmc}): the projection has been undone.
Computed by \texttt{manybodymc.py}.}
\label{tab:16-spinmoves}
\end{table}

The remedy is to exchange spins during the variational equilibration, which
leaves each walker in a randomised assignment, and then freeze them.  The
ensemble still averages over assignments -- different walkers carry different
ones -- and nothing interferes with the projection.  That is the default in
\texttt{manybodymc.py}, and Table~\ref{tab:16-spinmoves} shows it agreeing with
the fully frozen calculation.

\paragraph{The time step.}
Table~\ref{tab:16-timestep} repeats the extrapolation of
Section~\ref{sec:dmctimestep}, with the equilibration measured in imaginary
time rather than in steps -- at $\Delta\tau=0.01$ a fixed number of steps
covers a tenth of the imaginary time it does at $\Delta\tau=0.1$, and a
calculation that forgets this reports the time-step error of its own
unconverged projection.

\begin{table}[h]
\centering
\renewcommand{\arraystretch}{1.2}
\setlength{\tabcolsep}{13pt}
\begin{tabular}{lll}
\toprule
$\Delta\tau$ & energy & acceptance\\
\midrule
$0.10$ & $20.16003(230)$ & $94.3\%$\\
$0.05$ & $20.16028(353)$ & $97.4\%$\\
$0.02$ & $20.15882(327)$ & $99.2\%$\\
$0.01$ & $20.16832(363)$ & $99.7\%$\\
\midrule
weighted mean & $20.16125(151)$ & $\chi^2/\mathrm{dof}=1.6$\\
\bottomrule
\end{tabular}
\caption{Diffusion Monte Carlo for six electrons at $\alpha=0.98$,
$\beta=0.46$, with $250$ walkers, $500$ steps and an equilibration of five
units of imaginary time at every time step.  No time-step dependence is
resolvable.  Computed by \texttt{manybodymc.py}.}
\label{tab:16-timestep}
\end{table}

No time-step dependence survives the error bars, so the weighted mean is the
result:
\begin{equation}
  \boxed{\;E_{\mathrm{DMC}} = 20.16125 \pm 0.00151\;}
  \label{eq:16-dmc}
\end{equation}
against the published $20.1597(2)$ of Eq.~(\ref{eq:16-benchmark}).  The two
agree to one standard error.

%----------------------------------------------------------------------
\section{The comparison}
\label{sec:appcomparison}

\begin{table}[h]
\centering
\renewcommand{\arraystretch}{1.25}
\setlength{\tabcolsep}{10pt}
\begin{tabular}{llll}
\toprule
method & energy & basis & where\\
\midrule
Hartree-Fock & $20.72026$ & $42$ orbitals
  & Table~\ref{tab:11-hierarchy}\\
MP2 & $20.30256$ & $42$ orbitals & Table~\ref{tab:11-hierarchy}\\
CCSD & $20.27325$ & $42$ orbitals & Table~\ref{tab:11-hierarchy}\\
CCSD & $20.1737$ & $20$ shells & Ref.~\cite{lohne2011}\\
$\Lambda$-CCSD(T) & $20.1582$ & $20$ shells & Ref.~\cite{lohne2011}\\
\midrule
VMC & $20.19961(91)$ & none & Eq.~(\ref{eq:16-vmc})\\
DMC & $20.16125(151)$ & none & Eq.~(\ref{eq:16-dmc})\\
DMC, published & $20.1597(2)$ & none & Ref.~\cite{lohne2011}\\
\bottomrule
\end{tabular}
\caption{Six electrons in a two-dimensional parabolic trap at
$\hbar\omega=1$, in atomic units.  The published coupled-cluster results use an
effective interaction and a self-consistent Hartree-Fock basis, so they are not
the same calculation as the $42$-orbital rows, only the same method carried to
convergence.}
\label{tab:16-comparison}
\end{table}

Three things can be read off Table~\ref{tab:16-comparison}, and they are
different in kind.

\emph{The basis is the limitation, not the method.}  CCSD in $42$ orbitals is
$0.11$ above CCSD in twenty shells, and almost the whole of the apparent gap
between coupled cluster and Monte Carlo is that, and not a failure of the
cluster expansion.  Read the other way: at convergence, $\Lambda$-CCSD(T) and
diffusion Monte Carlo -- a deterministic expansion in particle-hole excitations
and a stochastic projection with a fixed nodal surface, sharing no
approximation whatever -- agree to $1.5\times10^{-3}$ on a system with no exact
solution.  That agreement is the real result of this chapter, and neither
method could establish it alone.

\emph{The Monte Carlo hierarchy behaves as advertised.}  The variational energy
is an upper bound, the diffusion energy improves on it by $0.038$, and the
diffusion energy is itself an upper bound because the nodes are approximate.
Both are basis-free and neither has anything to converge.

\emph{The remaining error is nodal.}  Our $20.16125(151)$ and the published
$20.1597(2)$ use the same single-determinant nodes and agree; the difference
from the truth, whatever it is, is the fixed-node error, and no amount of
computer time reaches it.  Only better nodes will.

%----------------------------------------------------------------------
\section{Dependence on the confinement}
\label{sec:appomega}

Table~\ref{tab:11-omega} varied the trap frequency and found the correlation
energy growing as the trap shallows.  Table~\ref{tab:16-omega} adds the two
Monte Carlo columns, at $\alpha=0.98$ and $\beta=0.46\sqrt{\omega}$ -- the
oscillator length is $\omega^{-1/2}$, so that is the only scaling the Jastrow
range needs, and a short scan at each frequency confirms it.

\begin{table}[h]
\centering
\renewcommand{\arraystretch}{1.2}
\setlength{\tabcolsep}{7pt}
\small
\begin{tabular}{rlllll}
\toprule
$\hbar\omega$ & $E_{\mathrm{HF}}$ & $E_{\mathrm{MP2}}$ & $E_{\mathrm{CCSD}}$
  & $E_{\mathrm{VMC}}$ & $E_{\mathrm{DMC}}$\\
\midrule
$0.25$ & $\phantom{0}7.39179$ & $\phantom{0}7.06226$ & $\phantom{0}7.03671$
  & $\phantom{0}7.02274(89)$ & $\phantom{0}6.98238(106)$\\
$0.50$ & $12.27150$ & $11.89288$ & $11.86345$
  & $11.82437(199)$ & $11.78488(116)$\\
$1.00$ & $20.72026$ & $20.30256$ & $20.27325$
  & $20.19961(91)$ & $20.15960(278)$\\
$2.00$ & $35.67233$ & $35.22533$ & $35.19959$
  & $35.09581(341)$ & $35.04162(394)$\\
$4.00$ & $62.74116$ & $62.27348$ & $62.25276$
  & $62.11527(657)$ & $62.05754(472)$\\
\bottomrule
\end{tabular}
\caption{Six electrons at five trap frequencies.  The first three columns are
the $42$-orbital results of Table~\ref{tab:11-omega}; the last two are this
chapter, with $600$ walkers and $\Delta\tau=0.05/\omega$.  Computed by
\texttt{manybodymc.py}.}
\label{tab:16-omega}
\end{table}

The ordering
$E_{\mathrm{HF}}>E_{\mathrm{MP2}}>E_{\mathrm{CCSD}}>E_{\mathrm{VMC}}>
E_{\mathrm{DMC}}$ holds at every frequency, and the gap between the
coupled-cluster and the Monte Carlo columns widens as $\omega$ falls --
$0.014$ at $\omega=4$ against $0.054$ at $\omega=0.25$ -- because a shallow
trap is a more strongly correlated system and the finite basis is worse there.

\begin{notebox}
\textbf{An honest discrepancy.}  At $\hbar\omega=0.5$ the published diffusion
Monte Carlo value is $11.7888(2)$~\cite{lohne2011} and the table gives
$11.78488(116)$, which is below it -- impossible for two calculations sharing
the same nodes, since fixed-node energies are upper bounds.  The culprit is
population control.  Repeating that point with $250$ walkers instead of $600$
gives $11.77981(144)$, and with $800$ it gives $11.78543(149)$: the energy
climbs as the population grows, which is the standard signature of
population-control bias and the standard direction for it.  At $\hbar\omega=1$
the same test finds no movement between $100$ and $400$ walkers, because the
local energy fluctuates less and the branching is gentler.  The bias is
$\mathcal{O}(1/N_{\mathrm{walkers}})$ and would be removed by running larger
populations; it is left visible here because a residual of $4\times10^{-3}$ at
the shallowest traps is exactly the kind of thing a table of six-digit numbers
is apt to conceal.
\end{notebox}

%----------------------------------------------------------------------
\section{What each method costs}
\label{sec:appcost}

\begin{table}[h]
\centering
\renewcommand{\arraystretch}{1.25}
\setlength{\tabcolsep}{9pt}
\begin{tabular}{lll}
\toprule
method & scaling & limited by\\
\midrule
Hartree-Fock & $\bigO(n^4)$ per iteration & the basis\\
MP2 & $\bigO(n^5)$ & the basis\\
CCSD & $\bigO(n^6)$ per iteration & the basis\\
$\Lambda$-CCSD(T) & $\bigO(n^7)$ & the basis, weakly\\
\midrule
VMC & $\bigO(N^2)$ per sweep, $\bigO(M^{-1/2})$ & the trial function\\
DMC & as VMC & the nodes of the trial function\\
\bottomrule
\end{tabular}
\caption{$n$ is the number of single-particle states, $N$ the number of
electrons, $M$ the number of Monte Carlo samples.  The two families are limited
by entirely different things, which is why running both is worth more than
running either twice.}
\label{tab:16-cost}
\end{table}

The distinction in the last column is the point.  A basis-set method has a
systematic error that shrinks as $n$ grows and a statistical error of zero; a
Monte Carlo method has no basis-set error at all, a statistical error that
shrinks as $M^{-1/2}$, and -- for diffusion Monte Carlo -- one systematic error
that no amount of $M$ will touch.  When the two agree, as they do at the bottom
of Table~\ref{tab:16-comparison}, the agreement means something precisely
because the failure modes do not overlap.

The costs also scale in the number of \emph{particles} quite differently.
Going from six electrons to twenty multiplies the coupled-cluster work by
$(n_{20}/n_6)^6$ with a basis that must also grow, while the Monte Carlo work
grows as $N^2$ per sweep with the number of samples fixed by the accuracy
wanted.  That is why the largest circular dots in the literature are done by
diffusion Monte Carlo, and why Ref.~\cite{lohne2011} needed a renormalised
interaction to bring coupled cluster into agreement with it at twenty
electrons.

%----------------------------------------------------------------------
\section{Summary}
\label{sec:appsummary}

Six electrons is where the simplifications end.  The trial function needs a
determinant, and therefore the inverse Slater matrix of
Section~\ref{sec:slaterupdate}, the derivative ratios of
Section~\ref{sec:slatergradients} and the spin factorisation of
Section~\ref{sec:spinfactorisation}; the cusp coefficients become spin
dependent, $1$ for antiparallel pairs and $\tfrac13$ for parallel ones, because
the determinant supplies part of the node itself; the wave function acquires a
nodal surface, which makes the drift divergent and requires either Umrigar's
cutoff or the spin-exchange move to control; and diffusion Monte Carlo acquires
a fixed-node error that is now the accuracy limit.

Sampling the spin assignment is exact and useful in variational Monte Carlo,
where it confirms the Moskowitz-Kalos factorisation numerically and rescues
walkers from nodal regions, and it is wrong in diffusion Monte Carlo, where it
preserves $|\Psi_T|^2$ rather than $\Psi_T\Phi$ and undoes the projection
almost entirely.  The same move, correct in one method and destructive in the
other, is a compact reminder of what the two methods actually sample.

The numbers, finally, hang together.  A variational calculation with two
parameters gives $20.19961(91)$; projecting it gives $20.16125(151)$, in
agreement with published diffusion Monte Carlo at one standard error and with
converged $\Lambda$-CCSD(T) at $1.5\times10^{-3}$.  The $42$-orbital coupled
cluster of Chapter~\ref{chap:applications} sits a tenth of an atomic unit
above, and that gap is the basis and nothing else.

What no method here can do is improve the nodes.  Both Monte Carlo results rest
on a single Slater determinant of oscillator orbitals, and
Eq.~(\ref{eq:16-node}) -- three electrons collinear -- is a guess about the
ground state that was never optimised and, within this ansatz, cannot be.  That
is the last of the approximations in this book to remain uncontrolled, and it
is the one that motivates what follows: trial functions flexible enough that
their nodes are learned rather than assumed.

%=============================================================================
\section{Exercises}
%=============================================================================

\subsection*{Exercise 1: the trial function}
\label{ex:16-trial}
\addcontentsline{toc}{subsection}{Exercise 1: the trial function}

\begin{enumerate}
\item[a)] Derive Eq.~(\ref{eq:16-cusp}) by requiring the local energy of a
  relative wave function behaving as $r^{\ell}$ at contact to stay finite as
  $r\to0$ in $d$ dimensions.  Evaluate it for $\ell=0,1$ in $d=2$ and $d=3$ and
  check the two-dimensional values against Eq.~(\ref{eq:16-cuspvalues}).

\item[b)] Verify by hand that Eq.~(\ref{eq:16-slaterjastrow}) reduces to the
  two-electron trial function of Eq.~(\ref{eq:13-trial}) when $N=2$, and say
  which of the two cusp coefficients survives and why.

\item[c)] Derive Eq.~(\ref{eq:16-laplacian}) and explain why the determinant
  and the Jastrow factor never have to be differentiated together.

\item[d)] The literature~\cite{harju2002} uses two range parameters, one for
  parallel and one for antiparallel pairs.  Implement this in
  \texttt{manybodymc.py} and optimise both with the gradient of
  Eq.~(\ref{eq:14-gradient}).  How much does the variational energy improve,
  and does the diffusion Monte Carlo energy change at all?
\end{enumerate}

\subsection*{Exercise 2: determinants in a Monte Carlo loop}
\label{ex:16-determinant}
\addcontentsline{toc}{subsection}{Exercise 2: determinants in a Monte Carlo loop}

\begin{enumerate}
\item[a)] Starting from Eq.~(\ref{eq:2-Rinverse}), show that when the
  determinant factorises as in Eq.~(\ref{eq:2-detfactorises}) a single-particle
  move leaves one factor untouched, and count the operations saved.

\item[b)] Replace the Sherman-Morrison update in \texttt{manybodymc.py} by a
  full recomputation of the inverse at every accepted move and measure the
  running time for $N=6$, $N=12$ and $N=20$.  At which $N$ does the asymptotic
  argument of Table~\ref{tab:vmcscaling} start to pay?

\item[c)] Monitor $\|\bm{D}^{-1}\bm{D}-\bm{1}\|$ over a long run.  Under what
  circumstances does it grow, and what does
  Section~\ref{sec:slaterstability} suggest doing about it?

\item[d)] Show that for six electrons the nodal surface of
  $\Det{D}_{\uparrow}$ is the collinearity condition of
  Eq.~(\ref{eq:16-node}), and that it is unchanged by the scale parameter
  $\alpha$.  What is the corresponding condition for twelve electrons?
\end{enumerate}

\subsection*{Exercise 3: sampling the spins}
\label{ex:16-spin}
\addcontentsline{toc}{subsection}{Exercise 3: sampling the spins}

\begin{enumerate}
\item[a)] Show that the exchange proposal of Section~\ref{sec:appspin} is
  symmetric, so that Eq.~(\ref{eq:16-spinmetropolis}) is the correct
  acceptance ratio, and that it conserves $S_z$.

\item[b)] Set $a_{\uparrow\uparrow}=a_{\uparrow\downarrow}$ in
  \texttt{manybodymc.py} and measure the exchange acceptance rate.  Explain the
  result in terms of what the move can and cannot change.

\item[c)] Reproduce Table~\ref{tab:16-nodes}.  Then record, over a long run,
  the joint distribution of the local energy and $\ln|\Det{D}|$, and use it to
  argue that the walkers with extreme local energies are those close to a node.

\item[d)] Reproduce Table~\ref{tab:16-spinmoves} and explain, from
  Eq.~(\ref{eq:15-importanceequation}), why a move that satisfies detailed
  balance with respect to $|\Psi_T|^2$ biases a diffusion Monte Carlo
  calculation.  Devise a numerical test that measures the size of the bias as a
  function of how often the exchange is attempted.
\end{enumerate}

\subsection*{Exercise 4: the comparison}
\label{ex:16-comparison}
\addcontentsline{toc}{subsection}{Exercise 4: the comparison}

\begin{enumerate}
\item[a)] Using \texttt{quantumdot.py}, extend the Hartree-Fock and CCSD
  calculations of Chapter~\ref{chap:applications} to as many shells as your
  machine allows, and plot the energy against the number of shells.  Fit the
  tail and extrapolate.  How close do you get to the $20.1737$ of
  Ref.~\cite{lohne2011}?

\item[b)] Repeat the population-size study of the note in
  Section~\ref{sec:appomega} at $\hbar\omega=0.25$, and estimate the
  population-control bias by extrapolating in
  $1/N_{\mathrm{walkers}}$.

\item[c)] Run both Monte Carlo methods for $N=12$, which fills three shells.
  Compare with the coupled-cluster results of Ref.~\cite{lohne2011}.  What
  changes in the cost, and what changes in the quality of the trial function?

\item[d)] Table~\ref{tab:16-cost} claims the two families of method are limited
  by different things.  Design a calculation that would distinguish a
  basis-set error from a fixed-node error in a system where neither is known
  independently, and say what it would cost.
\end{enumerate}
